\section*{Reading notes}
\subsection*{\texttt{GELFAND-P001455}}
\par\noindent\leavevmode\phantomsection\label{GELFAND-0158-U145501}\label{GELFAND-P001455:configuration-cancel}\textbf{\texttt{GELFAND-0158-U145501}} (configuration-cancel).\nopagebreak[4]\par
Between “Fuchs” and “not”, one sign is densely crossed out; an upper sloping stroke and a lower return are visible.\par
Candidates: у; д. Unresolved in the diplomatic source.\par
Historical source review: Original coordinates [1310,827,1520,930], magnification 7×; also the line [940,730,2030,960], 3×. The crossed-out beginning was not reconstructed from its meaning.\par
Current source comparison: The cancellation still prevents a secure letter reading. Neither candidate is promoted to a reading.\par\medskip
\par\noindent\leavevmode\phantomsection\label{GELFAND-0158-U145502}\label{GELFAND-P001455:task-number}\textbf{\texttt{GELFAND-0158-U145502}} (task-number).\nopagebreak[4]\par
After the underlined “Problem” is a sign with a loop and a long upper stroke, followed by a full stop.\par
Candidates: 6; Cyrillic в. Unresolved in the diplomatic source.\par
Historical source review: Original coordinates [480,1100,600,1210], 6×; line [200,1100,2150,1510], 3×. The question of numeral versus letter was left open.\par
Current source comparison: The loop and long rising stroke remain compatible with both the numeral and the Cyrillic letter.\par\medskip
\par\noindent\leavevmode\phantomsection\label{GELFAND-0158-U145503}\label{GELFAND-P001455:vertex-count-symbol}\textbf{\texttt{GELFAND-0158-U145503}} (vertex-count-symbol).\nopagebreak[4]\par
In the definition of the number of vertices and formulas (13)–(16′), a sign resembling handwritten ν and v recurs, with subscript i or a numeral.\par
Candidates: \textbackslash{}nu\_i; v\_i. Unresolved in the diplomatic source.\par
Historical source review: Checked [270,1515,2160,2190], 3×, and [115,2200,1730,3150], 3×. The representative ν was used in subsequent formulas; the choice of alphabet remains uncertain. The separate sign in the definition was additionally checked at [765,1935,910,2065], 10×; the ν/v distinction was not resolved.\par
Current source comparison: The mathematical role is clear, but it cannot decide the alphabet of this curved handwritten glyph.\par\medskip
\par\noindent\leavevmode\phantomsection\label{GELFAND-0158-U145504}\label{GELFAND-P001455:cancelled-equation-number}\textbf{\texttt{GELFAND-0158-U145504}} (cancelled-equation-number).\nopagebreak[4]\par
To the left of the last formula is a densely crossed-out parenthesized number; below and to the right, (16′) is legible.\par
Candidates: (17); (16). Unresolved in the diplomatic source.\par
Historical source review: Original coordinates [120,3290,355,3410], 7×; also [100,3190,1920,3440], 3×. The readable number (16′) was preserved independently of the crossed-out number.\par
Current source comparison: The readable (16′) is retained. The overwritten earlier number remains unresolved.\par\medskip
\subsection*{\texttt{GELFAND-P001456}}
\par\noindent\leavevmode\phantomsection\label{GELFAND-0158-U145601}\label{GELFAND-P001456:initial-charge-cancel}\textbf{\texttt{GELFAND-0158-U145601}} (initial-charge-cancel).\nopagebreak[4]\par
A short expression before 60(6−i) is crossed out; the final i can be distinguished, while several horizontal strokes obscure the first part.\par
Candidates: 60-i; 6-i. Unresolved in the diplomatic source.\par
Historical source review: Original coordinates [870,1280,1110,1375], 8×, and [75,1045,2230,1555], 3×. The legible replacement 60(6−i) was preserved separately.\par
Current source comparison: The earlier expression remains under heavy cancellation. Its replacement is clear; the charge calculation does not recover the cancelled text.\par\medskip
\par\noindent\leavevmode\phantomsection\label{GELFAND-0158-U145602}\label{GELFAND-P001456:impossible-cancel}\textbf{\texttt{GELFAND-0158-U145602}} (impossible-cancel).\nopagebreak[4]\par
After “impossible” and before “in view of”, one small sign or short beginning of a word with a long lower stroke is crossed out.\par
Candidates: none recorded. Unresolved in the diplomatic source.\par
Historical source review: Original coordinates [1160,1785,1235,1900], 8×; context [65,1565,2210,2150], 3×. No reliable reading was established beneath the cancellation.\par
Current source comparison: The lower return and crossing strokes do not establish a word or single sign.\par\medskip
\par\noindent\leavevmode\phantomsection\label{GELFAND-0158-U145603}\label{GELFAND-P001456:situation-one-cancel}\textbf{\texttt{GELFAND-0158-U145603}} (situation-one-cancel).\nopagebreak[4]\par
Between “0 or 1” and “vert. [abbr.]” is a crossed-out word beginning with р and containing the tall stroke of б.\par
Candidates: рёбр [source abbreviation]; рёбер [edges]. Unresolved in the diplomatic source.\par
Historical source review: Original coordinates [1015,2315,1185,2425], 8×; line [55,2170,2010,2580], 3×. Strokes obscure the end of the word.\par
Current source comparison: The visible beginning remains compatible with the two historical candidates; the ending is not established.\par\medskip
\par\noindent\leavevmode\phantomsection\label{GELFAND-0158-U145604}\label{GELFAND-P001456:cancelled-enumeration}\textbf{\texttt{GELFAND-0158-U145604}} (cancelled-enumeration).\nopagebreak[4]\par
Between the explanation of (23) and line (24), a separate sign with a loop-shaped top and a lower intersection is crossed out.\par
Candidates: 2. Unresolved in the diplomatic source.\par
Historical source review: Original coordinates [330,2580,410,2685], 8×. Possibly an unfinished numeral 2; its purpose was not asserted.\par
Current source comparison: The isolated cancellation does not establish either a numeral or its function.\par\medskip
\par\noindent\leavevmode\phantomsection\label{GELFAND-0158-U145605}\label{GELFAND-P001456:situation-two-cancel}\textbf{\texttt{GELFAND-0158-U145605}} (situation-two-cancel).\nopagebreak[4]\par
After “conn. [abbr.]”, a group beginning “ребр…” and ending with a separate с is crossed out.\par
Candidates: ребра с [edges/edge with; source wording]; ребро с [an edge with]. Unresolved in the diplomatic source.\par
Historical source review: Original coordinates [765,2700,1040,2820], 7×, and [65,2580,2070,2870], 3×. The subsequent legible phrase “with ≤2 vertices of valency 5” was left as in the source.\par
Current source comparison: The crossed-out ending remains unresolved. The following ≤2 was independently inspected at 3×/5×/10× and is literal; the separate discharging calculation does not change it.\par\medskip
\par\noindent\leavevmode\phantomsection\label{GELFAND-0158-U145606}\label{GELFAND-P001456:vertex-count-symbol}\textbf{\texttt{GELFAND-0158-U145606}} (vertex-count-symbol).\nopagebreak[4]\par
In (17) and (19), the handwritten sign for the number of vertices recurs and admits ν or v.\par
Candidates: \textbackslash{}nu; v. Unresolved in the diplomatic source.\par
Historical source review: The top of the page [70,30,1990,295], 3×, was compared with the definition on P001455. Here and subsequently, the formula uses ν as a representative without resolving the uncertainty about its alphabet.\par
Current source comparison: The same curved glyph recurs. The representative ν and its uncertainty are retained.\par\medskip
\par\noindent\leavevmode\phantomsection\label{GELFAND-0158-U145607}\label{GELFAND-P001456:diagram-18-right}\textbf{\texttt{GELFAND-0158-U145607}} (diagram-18-right).\nopagebreak[4]\par
In the right-hand pentagonal drawing in (18), the upper-right edge is emphasized with a dense wedge-shaped stroke; a protrusion appears near its right end. The remaining outline, central node and five radial segments can be distinguished.\par
Candidates: A thickened or retraced edge with a residual stroke; An edge bearing an additional arrow-like or other mark. Unresolved in the diplomatic source.\par
Historical source review: Original coordinates [80,300,1680,710], magnification 3×; the upper-right side is approximately [990,310,1110,390]. The nature of the emphasis was not inferred from graph theory. SourceFigure preserves a description of the entire third figure and its caption. Additionally checked [990,305,1115,405], 10×: the dense wedge and protrusion are visible, but their purpose remains uncertain.\par
Current source comparison: The wedge and short protrusion are traced as visible strokes, without asserting a directed edge. The partial recreation accompanies the exact original crop.\par\medskip
\subsection*{\texttt{GELFAND-P001457}}
\par\noindent\leavevmode\phantomsection\label{GELFAND-0158-U145701}\label{GELFAND-P001457:vanishing-cycles-cancel}\textbf{\texttt{GELFAND-0158-U145701}} (vanishing-cycles-cancel).\nopagebreak[4]\par
In the explanation to the right of formula (2), a word before the legible “cycles” is densely crossed out; several lower loops and a long right-hand descender can be distinguished beneath horizontal strokes.\par
Candidates: коциклов [cocycles]; когомол. [cohomol.; source abbreviation]. Unresolved in the diplomatic source.\par
Historical source review: Checked original coordinates [1715,2550,1970,2670], 8×, and the wider area [1700,2550,2040,2690], 6×; context [115,2380,2360,2840], 3×. Several strokes obscure the word, and the readings remain conjectural.\par
Current source comparison: The cancellation still does not distinguish the historical candidates.\par\medskip
\par\noindent\leavevmode\phantomsection\label{GELFAND-0158-U145702}\label{GELFAND-P001457:hypercohomology-word}\textbf{\texttt{GELFAND-0158-U145702}} (hypercohomology-word).\nopagebreak[4]\par
Below the condition ε≠0, a small parenthetical word begins “гипер…” and ends “…гомологии”; a separate “ко” between р and го cannot be reliably distinguished.\par
Candidates: гипергомологии [hyperhomology]; гиперкогомологии [hypercohomology]. Unresolved in the diplomatic source.\par
Historical source review: Original coordinates [430,2730,910,2845], magnification 6×; full context [115,2380,2360,2840], 3×. An independent repeated examination also did not resolve the doubt; the mathematically expected term was not substituted.\par
Current source comparison: The joined writing favors no certified insertion of ко. Both literal alternatives are preserved.\par\medskip
\subsection*{\texttt{GELFAND-P001458}}
\par\noindent\leavevmode\phantomsection\label{GELFAND-0158-U145801}\label{GELFAND-P001458:top-cancellation}\textbf{\texttt{GELFAND-0158-U145801}} (top-cancellation).\nopagebreak[4]\par
In the first line, between “that,” and “here”, a word is crossed out with horizontal strokes. “при чём” is clearly inserted above it. The original frame is 2408×3480; region (795,30,1130,150).\par
Candidates: почему [why]; зачем [for what purpose]. Unresolved in the diplomatic source.\par
Historical source review: The region was examined at 6×; the initial and middle letters are obscured. The candidates refer specifically to the crossed-out text and do not replace the insertion.\par
Current source comparison: The clear insertion and the cancelled earlier word remain distinct. No reading beneath the hatching is established.\par\medskip
\par\noindent\leavevmode\phantomsection\label{GELFAND-0158-U145802}\label{GELFAND-P001458:vanishing-cycle-sketch}\textbf{\texttt{GELFAND-0158-U145802}} (vanishing-cycle-sketch).\nopagebreak[4]\par
At the left below the first line, upper arcs converge on the left; a narrow hatched oval is on the right. Below is a horizontal oval with two points and a connecting stroke with a left-hand head. A broken line descends from the upper right-hand oval. Original region (285,145,1190,440).\par
Candidates: none recorded. Unresolved in the diplomatic source.\par
Historical source review: Examined at 3×. Intersections and thickening prevent a confident recreation of all directions and geometry; SourceFigure preserves the connection to the original drawing. The caption has a separate reading record.\par
Current source comparison: The hatched narrow oval and lower left-hand head are visible at 5×/10×. The partial recreation retains the broken vertical strokes without turning them into an arrow. The original crop remains alongside it.\par\medskip
\par\noindent\leavevmode\phantomsection\label{GELFAND-0158-U145803}\label{GELFAND-P001458:cycle-caption}\textbf{\texttt{GELFAND-0158-U145803}} (cycle-caption).\nopagebreak[4]\par
Caption to the right of the drawing: “цикл стя[н/г]. в точку” [cycle contracts to a point; uncertain source abbreviation]. The last letter of the abbreviation cannot be distinguished between н and г. Region (688,205,1155,313).\par
Candidates: стян. [source abbreviation]; стяг. [source abbreviation]. Unresolved in the diplomatic source.\par
Historical source review: Examined at 6×; an independent repeated examination also retained both alternatives, with a preference for стян. The abbreviation was not expanded.\par
Current source comparison: The ending remains a thick joined stroke. The recreation types both abbreviated alternatives, and the English text marks the source abbreviation rather than claiming an exact expansion.\par\medskip
\par\noindent\leavevmode\phantomsection\label{GELFAND-0158-U145804}\label{GELFAND-P001458:vanishing-cycles-degree}\textbf{\texttt{GELFAND-0158-U145804}} (vanishing-cycles-degree).\nopagebreak[4]\par
In the definition after “i.e.” is Φ with lower f and a small dark upper sign, then (𝓕)ₓ=0. The upper sign slopes and has a thickening and a small opening. Region (1987,439,2130,535).\par
Candidates: \textbackslash{}Phi\_f\textasciicircum{}i(\textbackslash{}mathcal F)\_x=0; \textbackslash{}Phi\_f\textasciicircum{}\textbackslash{}bullet(\textbackslash{}mathcal F)\_x=0; \textbackslash{}Phi\_f\textasciicircum{}*(\textbackslash{}mathcal F)\_x=0. Unresolved in the diplomatic source.\par
Historical source review: Magnification 10× did not distinguish the superscript. The context of the vanishing-cycle complex permits an index or grading sign, but mathematical expectation was not used to choose. An independent repeated examination also leaned toward i without removing the uncertainty.\par
Current source comparison: The sloping loop-like mark remains compatible with the historical alternatives; the draft does not resolve its degree.\par\medskip
\par\noindent\leavevmode\phantomsection\label{GELFAND-0158-U145805}\label{GELFAND-P001458:cotangent-cancellation}\textbf{\texttt{GELFAND-0158-U145805}} (cotangent-cancellation).\nopagebreak[4]\par
After df(X)⊂ and before Ω¹, several thick intersecting strokes obscure a cancelled sign. Region (1912,945,2160,1100).\par
Candidates: none recorded. Unresolved in the diplomatic source.\par
Historical source review: Examined at 6×; the cancelled symbol was not reconstructed from the expected notation for the cotangent bundle.\par
Current source comparison: The heavy intersections remain undeciphered. No cotangent notation is supplied beneath them.\par\medskip
\par\noindent\leavevmode\phantomsection\label{GELFAND-0158-U145806}\label{GELFAND-P001458:normalization-argument}\textbf{\texttt{GELFAND-0158-U145806}} (normalization-argument).\nopagebreak[4]\par
In the line “for a Morse f-n f with df(x)=ν”, f(0)=0 is then clearly written, with a round zero in the argument. Region (275,1750,2180,2160), especially the line around y=2040–2120.\par
Candidates: f(0)=0 — literal reading of the source; f(x)=0 — a possible mathematical normalization, not asserted in the source. Unresolved in the diplomatic source.\par
Historical source review: The literal f(0)=0 was preserved after 3× inspection. The preceding formula, however, specifies the fiber at x; perhaps x=0 is a local coordinate, or perhaps this is an error in the notes. This is a mathematical question, not an illegible numeral. No correction was made to the body.\par
Current source comparison: The first narrow 3× crop was blank and supplied no evidence. A wider 3× context and corrected 5× detail locate the formula; overlapping exact 10× details confirm df(x)=ν and f(0)=0. Choosing x as a local origin is a possible explanation, not a textual emendation.\par\medskip
\par\noindent\leavevmode\phantomsection\label{GELFAND-0158-U145807}\label{GELFAND-P001458:lower-detached-mark}\textbf{\texttt{GELFAND-0158-U145807}} (lower-detached-mark).\nopagebreak[4]\par
Above “if u∈𝓕\textasciicircum{}\{0,0\}” is a separate small sloping sign with two crossing strokes, resembling 4 or an insertion mark. Region (412,3090,800,3240).\par
Candidates: 4; insertion mark. Unresolved in the diplomatic source.\par
Historical source review: Examined at 6×. It was not treated as a page number or an established item number; its position above the beginning of the line was preserved.\par
Current source comparison: The detached cross-shaped mark still does not establish a numbered item.\par\medskip
\subsection*{\texttt{GELFAND-P001459}}
\par\noindent\leavevmode\phantomsection\label{GELFAND-0158-U145901}\label{GELFAND-P001459:upper-cancellation}\textbf{\texttt{GELFAND-0158-U145901}} (upper-cancellation).\nopagebreak[4]\par
Before “perverse sheaves [abbr.]” and after “are obtained”, a short word is crossed out several times. Region (1835,325,2030,392).\par
Candidates: none recorded. Unresolved in the diplomatic source.\par
Historical source review: Examined at 10×; the original letters have merged with the cancellation strokes. The presence of a cancelled word was preserved.\par
Current source comparison: No underlying word can be reliably separated from the horizontal cancellation.\par\medskip
\par\noindent\leavevmode\phantomsection\label{GELFAND-0158-U145902}\label{GELFAND-P001459:bottom-cutoff}\textbf{\texttt{GELFAND-0158-U145902}} (bottom-cutoff).\nopagebreak[4]\par
Below a long horizontal line at the bottom edge of the short text fragment, small upper portions of strokes remain; the left cluster resembles М. Region (510,657,1520,712).\par
Candidates: М…; illegible cropped strokes. Unresolved in the diplomatic source.\par
Historical source review: Examined at 6×. No continuation was invented; all visible cropped traces were accounted for. The rest of the frame below this fragment is blank.\par
Current source comparison: The enlarged left cluster does not provide a complete letter or continuation. The cutoff remains explicit.\par\medskip
\subsection*{\texttt{GELFAND-P001460}}
\par\noindent\leavevmode\phantomsection\label{GELFAND-0158-U146002}\label{GELFAND-P001460:stable-cancellation}\textbf{\texttt{GELFAND-0158-U146002}} (stable-cancellation).\nopagebreak[4]\par
In the line “Stable sheaves [abbr.] and”, the following word is crossed out with several thick horizontal strokes; no reliably complete letters survive. Region (870,2890,1080,2960).\par
Candidates: none recorded. Unresolved in the diplomatic source.\par
Historical source review: Examined at 10×, and also the line (445,2870,1090,2980) at 3×. There is no clear closing quotation mark after the cancellation; none was added.\par
Current source comparison: The underlying word and a closing quotation mark are not established.\par\medskip
\par\noindent\leavevmode\phantomsection\label{GELFAND-0158-U146003}\label{GELFAND-P001460:quantum-title-mark}\textbf{\texttt{GELFAND-0158-U146003}} (quantum-title-mark).\nopagebreak[4]\par
Between “Quantum groups [abbr.]” and “(after Drinfeld)” is a small dense sign above the line: upper horizontal strokes and two lower protrusions. Region (830,2985,928,3055).\par
Candidates: crossed-out closing quotation marks; asterisk. Unresolved in the diplomatic source.\par
Historical source review: Examined at 10×. The sign resembles filled-in quotation marks but was not replaced with a certain quotation mark or footnote sign. Separate closing quotation marks occur after the parentheses in the source.\par
Current source comparison: Both historical interpretations remain possible; the separate quotation marks after the parentheses are retained independently.\par\medskip
\subsection*{\texttt{GELFAND-P001461}}
\par\noindent\leavevmode\phantomsection\label{GELFAND-0158-U146101}\label{GELFAND-P001461:surface-sketch}\textbf{\texttt{GELFAND-0158-U146101}} (surface-sketch).\nopagebreak[4]\par
Drawing on the left in the middle of the page, original coordinates (170,1735,525,1985): an inner oval outline with two small loop-shaped strokes, and a broad outer outline with broken portions. The caption “рим. пов. ⊂ сфера” [Riemann surface ⊂ sphere; source abbreviations] is legible and is separately transcribed in the text.\par
Candidates: Two internal loop-shaped marks; the connection of the upper one to the outline is not established; A schematic depiction of handles of a surface. Unresolved in the diplomatic source.\par
Historical source review: The original 2408×3480 image and the 3× enlarged crop (170,1240,1310,2075) were inspected. The exact topology, the upper loop connection and the differing line types could not be confidently reconstructed; SourceFigure preserves the original reference without inventing topology.\par
Current source comparison: At 5×/10×, the upper loop stops short of the oval outline and the lower mark has an open stroke and tail. The partial recreation preserves these gaps and the separate outer strokes; it asserts no genus or handle attachment.\par\medskip
\subsection*{\texttt{GELFAND-P001462}}
\par\noindent\leavevmode\phantomsection\label{GELFAND-0158-U146201}\label{GELFAND-P001462:antipode-cancelled-sign}\textbf{\texttt{GELFAND-0158-U146201}} (antipode-cancelled-sign).\nopagebreak[4]\par
After item number “3.” and before “S:A→A”, a single densely crossed-out sign appears. Original coordinates (345,1260,430,1388).\par
Candidates: Δ; Cyrillic и. Unresolved in the diplomatic source.\par
Historical source review: The original and a 10× crop were examined. Several diagonal strokes obscure the sign; neither candidate was asserted. The crossed-out area was not omitted. An independent repeated examination by the root reviewer at 10× also failed to distinguish the sign unambiguously.\par
Current source comparison: The repeated current 3×/5×/10× inspection leaves the sign unresolved. The root-review claim in the translated historical note describes that earlier review.\par\medskip
\par\noindent\leavevmode\phantomsection\label{GELFAND-0158-U146202}\label{GELFAND-P001462:antipode-formula}\textbf{\texttt{GELFAND-0158-U146202}} (antipode-formula).\nopagebreak[4]\par
Formula (4) is explicitly written S(X\textasciicircum{}±)=q\textasciicircum{}\{H/2\}X\textasciicircum{}±q\textasciicircum{}\{-H/2\}, without an initial minus on the right-hand side. Original crop (410,1390,1290,1610), magnification 5×.\par
Candidates: S(X\textasciicircum{}±)=q\textasciicircum{}\{H/2\}X\textasciicircum{}±q\textasciicircum{}\{-H/2\} — literal source formula. Unresolved in the diplomatic source.\par
Historical source review: A mathematical comparison with the antipode diagram shown below and with S(H)=−H reveals a possible sign or normalization error in the source. This is not uncertainty about a visible minus: there is none before q. The formula was preserved without correction.\par
Current source comparison: The absence of the initial minus is confirmed at 3×/5×/10×. The separate exact matrix calculation checks compatibility with the written coproduct and antipode axiom; it cannot change this literal formula.\par\medskip
\par\noindent\leavevmode\phantomsection\label{GELFAND-0158-U146203}\label{GELFAND-P001462:cancelled-generators}\textbf{\texttt{GELFAND-0158-U146203}} (cancelled-generators).\nopagebreak[4]\par
After “generators”, a fragment is densely crossed out; forms resembling K, X\textasciicircum{}± can be distinguished, with obscured loop-shaped strokes remaining on the right. Original coordinates (1530,2410,1915,2580), review crop (1340,2370,1925,2610), 5×.\par
Candidates: K,X\textasciicircum{}±, then illegible material; K,X\textasciicircum{}±, q…. Unresolved in the diplomatic source.\par
Historical source review: The later line “K=q\textasciicircum{}\{H/2\}, K,X\textasciicircum{}±” and the lower crossed-out K,X\textasciicircum{}± are separately legible and were not used to reconstruct this cancelled fragment.\par
Current source comparison: The readable later generator list does not decide the overwritten right-hand loops in this fragment.\par\medskip
\par\noindent\leavevmode\phantomsection\label{GELFAND-0158-U146204}\label{GELFAND-P001462:commutator-normalization}\textbf{\texttt{GELFAND-0158-U146204}} (commutator-normalization).\nopagebreak[4]\par
At the bottom is X\textasciicircum{}+X\textasciicircum{}−−X\textasciicircum{}−X\textasciicircum{}+=(K−K\textasciicircum{}\{−1\})/(q−q\textasciicircum{}\{−1\}), with K=q\textasciicircum{}\{H/2\} above. The preceding page GELFAND-P001461 ends with [X\textasciicircum{}+,X\textasciicircum{}−]=sh((h/2)H)/sh(h/2), q=e\textasciicircum{}h. Original crop of the lower formula (370,2925,1310,3180), 5×.\par
Candidates: (K−K\textasciicircum{}\{−1\})/(q−q\textasciicircum{}\{−1\}) — literal source formula. Unresolved in the diplomatic source.\par
Historical source review: Substituting the preceding notation shows that the denominators of the two visible formulas differ: sh(h/2) gives q\textasciicircum{}\{1/2\}−q\textasciicircum{}\{−1/2\}. No change of convention is indicated within these lines. The source was reproduced without correcting the denominator.\par
Current source comparison: The denominator q−q\textasciicircum{}\{-1\} is clear at 5×/10×. The exact comparison keeps the generators fixed and records the normalization discrepancy separately.\par\medskip
\subsection*{\texttt{GELFAND-P001463}}
\par\noindent\leavevmode\phantomsection\label{GELFAND-0158-U146301}\label{GELFAND-P001463:r-coefficient}\textbf{\texttt{GELFAND-0158-U146301}} (r-coefficient).\nopagebreak[4]\par
Formula (8), coefficient numerator: q\textasciicircum{}n; denominator [n]\_q!; followed by e\textasciicircum{}n⊗f\textasciicircum{}n. Coordinates in the original PNG: x=1040–1230, y=1160–1400; examined at 3× and 6×.\par
Candidates: q\textasciicircum{}n/[n]\_q!. Unresolved in the diplomatic source.\par
Historical source review: The exponent n is directly readable. No additional powers or factors are visible in this formula, and none were added. The mathematical completeness of the handwritten coefficient was left for checking.\par
Current source comparison: The numerator q\textasciicircum{}n is confirmed by the current enlargements. Any mathematical test must state its convention for [n]\_q; no familiar replacement coefficient is inserted.\par\medskip
\par\noindent\leavevmode\phantomsection\label{GELFAND-0158-U146302}\label{GELFAND-P001463:cancelled-before-membership}\textbf{\texttt{GELFAND-0158-U146302}} (cancelled-before-membership).\nopagebreak[4]\par
Before the parenthesis (∈ U\_q(SU(2))\textasciicircum{}\{⊗3\}), to the right of the first line of (10), several signs are densely crossed out. Original coordinates: x=1410–1525, y=1900–2045; examined at 3× and 10×.\par
Candidates: none recorded. Unresolved in the diplomatic source.\par
Historical source review: The cancellation strokes prevent reliable separation of the original letters; the cancellation was marked in the text and its content was not reconstructed.\par
Current source comparison: The diagonals and loops do not establish the cancelled signs.\par\medskip
\par\noindent\leavevmode\phantomsection\label{GELFAND-0158-U146303}\label{GELFAND-P001463:exercise-harmonics}\textbf{\texttt{GELFAND-0158-U146303}} (exercise-harmonics).\nopagebreak[4]\par
In the added note after “in spherical [abbr.]”, an abbreviation “гар-…м” appears: a small loop-shaped stroke is visible between the hyphen and final м. Original coordinates: x=1890–2175, y=1030–1105; examined at 3× and independently at 8×.\par
Candidates: гар-ям [source abbreviation]; гар-м [source abbreviation]. Unresolved in the diplomatic source.\par
Historical source review: The main line and the note “start with SU(2)” were read; only the abbreviation remains disputed. The abbreviation was not expanded.\par
Current source comparison: The small intermediate stroke remains ambiguous. The translation identifies the abbreviation and its alternatives rather than certifying an expansion.\par\medskip
\subsection*{\texttt{GELFAND-P001464}}
\par\noindent\leavevmode\phantomsection\label{GELFAND-0158-U146401}\label{GELFAND-P001464:definition-separator}\textbf{\texttt{GELFAND-0158-U146401}} (definition-separator).\nopagebreak[4]\par
In the first line, between R\textasciicircum{}\{j₁j₂\} and V\textasciicircum{}\{j₁\}, are two thickened elements, one above the other, with additional horizontal strokes. Original coordinates: x=490–567, y=142–235; examined at 3× and 10×.\par
Candidates: colon; colon over a corrected sign. Unresolved in the diplomatic source.\par
Historical source review: The formula was given with the candidate “:” and an explicit note. The original corrected sign was not established.\par
Current source comparison: The stacked thick marks support the displayed candidate colon but do not identify the previous sign.\par\medskip
\par\noindent\leavevmode\phantomsection\label{GELFAND-0158-U146402}\label{GELFAND-P001464:cancelled-coefficient}\textbf{\texttt{GELFAND-0158-U146402}} (cancelled-coefficient).\nopagebreak[4]\par
After “q-analogues” and before “Clebsch–Gordan coeff. [abbr.]”, one or more letter-like strokes are densely crossed out. Original coordinates: x=677–746, y=1110–1205; examined at 3× and 10×.\par
Candidates: none recorded. Unresolved in the diplomatic source.\par
Historical source review: No reliable candidates for the original crossed-out sign were obtained; it was not replaced with an expected word.\par
Current source comparison: The overwritten strokes remain undeciphered.\par\medskip
\par\noindent\leavevmode\phantomsection\label{GELFAND-0158-U146403}\label{GELFAND-P001464:sum-mark}\textbf{\texttt{GELFAND-0158-U146403}} (sum-mark).\nopagebreak[4]\par
In (15), above the sum on the right, a small circle has thickened vertical and horizontal lines and a diagonal stroke in its lower-left part. Original coordinates: x=1015–1085, y=1245–1310; examined at 3× and 10×.\par
Candidates: ⊕; ⊗ corrected to ⊕. Unresolved in the diplomatic source.\par
Historical source review: The formula preserves the legible final plus in a circle; the possible correction of an earlier sign remains unresolved. The summation sign was not replaced by a single direct-sum symbol.\par
Current source comparison: The final circled plus remains legible, with a possible earlier diagonal. Its correction history remains uncertain.\par\medskip
\par\noindent\leavevmode\phantomsection\label{GELFAND-0158-U146404}\label{GELFAND-P001464:coefficient-range}\textbf{\texttt{GELFAND-0158-U146404}} (coefficient-range).\nopagebreak[4]\par
At the bottom left is A with upper indices m′₁…m′\_N and lower indices m₁…m\_L. Original coordinates: x=370–813, y=2860–3090; examined at 3× and 5×.\par
Candidates: A\textasciicircum{}\{m′₁…m′\_N\}\_\{m₁…m\_L\}. Unresolved in the diplomatic source.\par
Historical source review: In the map A written above, N occurs on the source factors and L on the target factors; the diagrams also have lower j\_N and upper j′\_L. In the coefficient, the final indices N/L occur in the opposite rows. The visible formula was preserved without reconciling the notation.\par
Current source comparison: Overlapping exact 10× details confirm the printed N/L ordering. A tensor-dimension check concerns compatibility only and does not reverse these visible indices.\par\medskip
\par\noindent\leavevmode\phantomsection\label{GELFAND-0158-U146405}\label{GELFAND-P001464:diagram-02}\textbf{\texttt{GELFAND-0158-U146405}} (diagram-02).\nopagebreak[4]\par
Lower right-hand diagram: a rectangle A with three vertical terminals above and below, without arrows. Above are j′₁,j′₂,j′\_L and, higher, m′₁,m′₂,…,m′ with a heavily corrected final index; below are j₁,j₂,j\_N and m₁,m₂,m\_N. Diagram x=1370–2035,y=2615–3445; disputed index x=1870–2010,y=2640–2765. Diagram examined at 3×, index at 10×.\par
Candidates: upper-right index L over N; upper-right index L over an illegible sign; upper-right index N with a correction. Unresolved in the diplomatic source.\par
Historical source review: The initial letter beneath the heavy correction in the upper-right caption is unclear. The entire diagram was therefore left as SourceFigure with a description of all visible terminals and labels. The first diagram (17), with unambiguous labels, was transcribed as inline TikZ.\par
Current source comparison: The current partial native comparison types all clear labels and brackets the corrected last upper index as [L/N ?]. It preserves six unarrowed terminals and does not infer the overwritten letter from the adjacent coefficient or map.\par\medskip
